
\documentclass{article}
\usepackage{colortbl}
\usepackage{makecell}
\usepackage{multirow}
\usepackage{supertabular}

\begin{document}

\newcounter{utterance}

\centering \large Interaction Transcript for game `referencegame', experiment `random\_grids', episode 3 with nemotron{-}3{-}ultra.
\vspace{24pt}

{ \footnotesize  \setcounter{utterance}{1}
\setlength{\tabcolsep}{0pt}
\begin{supertabular}{c@{$\;$}|p{.15\linewidth}@{}p{.15\linewidth}p{.15\linewidth}p{.15\linewidth}p{.15\linewidth}p{.15\linewidth}}
    \# & $\;$A & \multicolumn{4}{c}{Game Master} & $\;\:$B\\
    \hline

    \theutterance \stepcounter{utterance}  
    & & \multicolumn{4}{p{0.6\linewidth}}{
        \cellcolor[rgb]{0.9,0.9,0.9}{
            \makecell[{{p{\linewidth}}}]{
                \texttt{\tiny{[P1$\langle$GM]}}
                \texttt{Három rácsot kapsz, mindegyik 5×5{-}ös méretű.} \\
\texttt{A rácsok üres cellái "▢" jellel, a kitöltött cellái pedig "X" jellel vannak jelölve.} \\
\texttt{A feladatod, hogy olyan leírást adj, amely egyértelműen azonosítja a célnak megadott rácsot, és megkülönbözteti a két zavaró rácstól.} \\
\texttt{Az első rács a célrács, a következő kettő pedig a zavaró rácsok.} \\
\\ 
\texttt{Célrács:} \\
\\ 
\texttt{X □ X X □} \\
\texttt{□ X □ X X} \\
\texttt{X X □ □ X} \\
\texttt{□ X X □ X} \\
\texttt{X □ X □ □} \\
\\ 
\texttt{Zavaró rács 1:} \\
\\ 
\texttt{□ □ X X □} \\
\texttt{X □ □ X X} \\
\texttt{□ X X □ X} \\
\texttt{X X □ □ X} \\
\texttt{□ X X □ □} \\
\\ 
\texttt{Zavaró rács 2:} \\
\\ 
\texttt{X X □ X □} \\
\texttt{□ X X □ X} \\
\texttt{X □ □ X □} \\
\texttt{□ X □ □ X} \\
\texttt{X □ X □ X} \\
\\ 
\texttt{Instrukció: Írd le a célrácsot.} \\
\texttt{A válaszodat a „Kifejezés: ” címkével kezdd. Más szöveget ne adj meg.} \\
\\ 
\\ 
\\ 
\\ 
\\ 
\\ 
\\ 
\\ 
\\ 
            }
        }
    }
    & \\ \\

    \theutterance \stepcounter{utterance}  
    & \multicolumn{4}{p{0.6\linewidth}}{
        \cellcolor[rgb]{0.9,0.9,0.9}{
            \makecell[{{p{\linewidth}}}]{
                \texttt{\tiny{[P1$\rangle$GM]}}
                \texttt{Kifejezés: A célrács a jobb alsó sarokból balra haladva az első sorban (5. sor) „X □ X □ □” mintát mutat, míg a két zavaró rács azonos sorában „□ □ X □ X” (Zavaró 1) respective „X □ X □ X” (Zavaró 2) áll. Emellett a célrács bal felső sarkában (1. sor, 1. oszlop) „X” van, Zavaró 1{-}ben „□”, Zavaró 2{-}ben „X”; a célrács 2. sor 1. oszlopa „□”, Zavaró 1{-}ben „X”, Zavaró 2{-}ben „□”; a célrács 3. sor 1. oszlopa „X”, Zavaró 1{-}ben „□”, Zavaró 2{-}ben „X”. Ezek a különbségek egyértelműen megkülönböztetik a célrácsot a zavaróktól.} \\
            }
        }
    }
    & & \\ \\

    \theutterance \stepcounter{utterance}  
    & & & \multicolumn{2}{p{0.3\linewidth}}{
        \cellcolor[rgb]{0.9,0.9,0.9}{
            \makecell[{{p{\linewidth}}}]{
                \texttt{\tiny{[GM$|$GM]}}
                \texttt{Kifejezés: A célrács a jobb alsó sarokból balra haladva az első sorban (5. sor) „X □ X □ □” mintát mutat, míg a két zavaró rács azonos sorában „□ □ X □ X” (Zavaró 1) respective „X □ X □ X” (Zavaró 2) áll. Emellett a célrács bal felső sarkában (1. sor, 1. oszlop) „X” van, Zavaró 1{-}ben „□”, Zavaró 2{-}ben „X”; a célrács 2. sor 1. oszlopa „□”, Zavaró 1{-}ben „X”, Zavaró 2{-}ben „□”; a célrács 3. sor 1. oszlopa „X”, Zavaró 1{-}ben „□”, Zavaró 2{-}ben „X”. Ezek a különbségek egyértelműen megkülönböztetik a célrácsot a zavaróktól.} \\
            }
        }
    }
    & & \\ \\

    \theutterance \stepcounter{utterance}  
    & & \multicolumn{4}{p{0.6\linewidth}}{
        \cellcolor[rgb]{0.9,0.9,0.9}{
            \makecell[{{p{\linewidth}}}]{
                \texttt{\tiny{[GM$\rangle$P2]}}
                \texttt{Három rácsot kapsz, mindegyik 5×5{-}ös méretű.} \\
\texttt{A rácsok üres cellái "▢" jellel vannak jelölve, a kitöltött cellák pedig "X" jellel.} \\
\texttt{Egy leírást is kapsz, amely az egyik rácsot írja le.} \\
\texttt{A feladatod az, hogy kiválaszd azt a rácsot, amelyik megfelel a kapott leírásnak.} \\
\texttt{A válaszodban csak azt add meg szövegesen, hogy az első, a második vagy a harmadik rácsról van{-}e szó.} \\
\\ 
\texttt{Első:} \\
\\ 
\texttt{X □ X X □} \\
\texttt{□ X □ X X} \\
\texttt{X X □ □ X} \\
\texttt{□ X X □ X} \\
\texttt{X □ X □ □} \\
\\ 
\texttt{Második:} \\
\\ 
\texttt{□ □ X X □} \\
\texttt{X □ □ X X} \\
\texttt{□ X X □ X} \\
\texttt{X X □ □ X} \\
\texttt{□ X X □ □} \\
\\ 
\texttt{Harmadik:} \\
\\ 
\texttt{X X □ X □} \\
\texttt{□ X X □ X} \\
\texttt{X □ □ X □} \\
\texttt{□ X □ □ X} \\
\texttt{X □ X □ X} \\
\\ 
\texttt{Kifejezés: A célrács a jobb alsó sarokból balra haladva az első sorban (5. sor) „X □ X □ □” mintát mutat, míg a két zavaró rács azonos sorában „□ □ X □ X” (Zavaró 1) respective „X □ X □ X” (Zavaró 2) áll. Emellett a célrács bal felső sarkában (1. sor, 1. oszlop) „X” van, Zavaró 1{-}ben „□”, Zavaró 2{-}ben „X”; a célrács 2. sor 1. oszlopa „□”, Zavaró 1{-}ben „X”, Zavaró 2{-}ben „□”; a célrács 3. sor 1. oszlopa „X”, Zavaró 1{-}ben „□”, Zavaró 2{-}ben „X”. Ezek a különbségek egyértelműen megkülönböztetik a célrácsot a zavaróktól.} \\
\texttt{Kérdés: Melyik rácsra vonatkozik a leírás?} \\
\texttt{A válaszodat a „Válasz: ” címkével kezdd, majd add meg a választ. Más szöveget ne adj meg.} \\
            }
        }
    }
    & \\ \\

    \theutterance \stepcounter{utterance}  
    & & & \multicolumn{4}{p{0.6\linewidth}}{
        \cellcolor[rgb]{0.9,0.9,0.9}{
            \makecell[{{p{\linewidth}}}]{
                \texttt{\tiny{[GM$\langle$P2]}}
                \texttt{Válasz: az első} \\
            }
        }
    }
     \\ \\

    \theutterance \stepcounter{utterance}  
    & & & \multicolumn{2}{p{0.3\linewidth}}{
        \cellcolor[rgb]{0.9,0.9,0.9}{
            \makecell[{{p{\linewidth}}}]{
                \texttt{\tiny{[GM$|$GM]}}
                \texttt{Invalid generated expression} \\
            }
        }
    }
    & & \\ \\

\end{supertabular}
}

\end{document}
